\section{Dimension of finite type algebras over fields}
\label{section-dimension-finite-type-algebras}

\noindent
In this section we compute the dimension of a polynomial ring over
a field. We also prove that the dimension of a finite type
domain over a field is the dimension of its local rings at maximal
ideals. We will establish the connection with the transcendence
degree over the ground field in
Section \ref{section-dimension-finite-type-algebras-reprise}.

\begin{lemma}
\label{lemma-dim-affine-space}
Let $\mathfrak m$ be a maximal ideal in $k[x_1, \ldots, x_n]$.
The ideal $\mathfrak m$ is generated by $n$ elements.
The dimension of $k[x_1, \ldots, x_n]_{\mathfrak m}$ is $n$.
Hence $k[x_1, \ldots, x_n]_{\mathfrak m}$ is a regular local
ring of dimension $n$.
\end{lemma}

\begin{proof}
By the Hilbert Nullstellensatz
(Theorem \ref{theorem-nullstellensatz})
we know the residue field $\kappa = \kappa(\mathfrak m)$ is
a finite extension of $k$. Denote $\alpha_i \in \kappa$ the
image of $x_i$. Denote $\kappa_i = k(\alpha_1, \ldots, \alpha_i)
\subset \kappa$, $i = 1, \ldots, n$ and $\kappa_0 = k$.
Note that $\kappa_i = k[\alpha_1, \ldots, \alpha_i]$
by field theory. Define inductively elements
$f_i \in \mathfrak m \cap k[x_1, \ldots, x_i]$
as follows: Let $P_i(T) \in \kappa_{i-1}[T]$
be the monic minimal polynomial of $\alpha_i $ over $\kappa_{i-1}$.
Let $Q_i(T) \in k[x_1, \ldots, x_{i-1}][T]$ be a monic lift of $P_i(T)$
(of the same degree). Set $f_i = Q_i(x_i)$.
Note that if $d_i = \deg_T(P_i) = \deg_T(Q_i) = \deg_{x_i}(f_i)$
then $d_1d_2\ldots d_i = [\kappa_i : k]$ by
Fields, Lemmas \ref{fields-lemma-multiplicativity-degrees} and
\ref{fields-lemma-degree-minimal-polynomial}.

\medskip\noindent
We claim that for all $i = 0, 1, \ldots, n$ there is an
isomorphism
$$
\psi_i : k[x_1, \ldots, x_i] /(f_1, \ldots, f_i) \cong \kappa_i.
$$
By construction the composition
$k[x_1, \ldots, x_i] \to k[x_1, \ldots, x_n] \to \kappa$
is surjective onto $\kappa_i$ and $f_1, \ldots, f_i$ are
in the kernel. This gives a surjective homomorphism.
We prove $\psi_i$ is injective by induction. It is clear for $i = 0$.
Given the statement for $i$ we prove it for $i + 1$.
The ring extension $k[x_1, \ldots, x_i]/(f_1, \ldots, f_i) \to
k[x_1, \ldots, x_{i + 1}]/(f_1, \ldots, f_{i + 1})$
is generated by $1$ element over a field and one
irreducible equation. By elementary field theory
$k[x_1, \ldots, x_{i + 1}]/(f_1, \ldots, f_{i + 1})$
is a field, and hence $\psi_i$ is injective.

\medskip\noindent
This implies that $\mathfrak m = (f_1, \ldots, f_n)$.
Moreover, we also conclude that
$$
k[x_1, \ldots, x_n]/(f_1, \ldots, f_i)
\cong
\kappa_i[x_{i + 1}, \ldots, x_n].
$$
Hence $(f_1, \ldots, f_i)$ is a prime ideal. Thus
$$
(0) \subset (f_1) \subset (f_1, f_2) \subset \ldots \subset
(f_1, \ldots, f_n) = \mathfrak m
$$
is a chain of primes of length $n$. The lemma follows.
\end{proof}

\begin{proposition}
\label{proposition-finite-gl-dim-polynomial-ring}
A polynomial algebra in $n$ variables over a field is a regular ring.
It has global dimension $n$. All localizations at maximal ideals
are regular local rings of dimension $n$.
\end{proposition}

\begin{proof}
By Lemma \ref{lemma-dim-affine-space}
all localizations $k[x_1, \ldots, x_n]_{\mathfrak m}$
at maximal ideals are regular local rings of dimension $n$. Hence
we conclude by Lemma \ref{lemma-finite-gl-dim-finite-dim-regular}.
\end{proof}

\begin{lemma}
\label{lemma-dimension-height-polynomial-ring}
Let $k$ be a field.
Let $\mathfrak p \subset \mathfrak q \subset k[x_1, \ldots, x_n]$
be a pair of primes.
Any maximal chain of primes between $\mathfrak p$ and $\mathfrak q$
has length $\text{height}(\mathfrak q) - \text{height}(\mathfrak p)$.
\end{lemma}

\begin{proof}
By
Proposition \ref{proposition-finite-gl-dim-polynomial-ring}
any local ring of $k[x_1, \ldots, x_n]$ is regular.
Hence all local rings are Cohen-Macaulay, see
Lemma \ref{lemma-regular-ring-CM}.
The local rings at maximal ideals have dimension $n$ hence
every maximal chain of primes in $k[x_1, \ldots, x_n]$
has length $n$, see
Lemma \ref{lemma-maximal-chain-CM}.
Hence every maximal chain of primes between $(0)$ and $\mathfrak p$
has length $\text{height}(\mathfrak p)$, see
Lemma \ref{lemma-CM-dim-formula}
for example.
Putting these together leads to the assertion of the lemma.
\end{proof}

\begin{lemma}
\label{lemma-dimension-spell-it-out}
Let $k$ be a field.
Let $S$ be a finite type $k$-algebra which is an integral domain.
Then $\dim(S) = \dim(S_{\mathfrak m})$ for any maximal
ideal $\mathfrak m$ of $S$. In words: every maximal chain
of primes has length equal to the dimension of $S$.
\end{lemma}

\begin{proof}
Write $S = k[x_1, \ldots, x_n]/\mathfrak p$.
By Proposition \ref{proposition-finite-gl-dim-polynomial-ring} and
Lemma \ref{lemma-dimension-height-polynomial-ring}
all the maximal chains of primes in $S$ (which necessarily end
with a maximal ideal) have length $n - \text{height}(\mathfrak p)$.
Thus this number is the dimension of $S$ and of $S_{\mathfrak m}$
for any maximal ideal $\mathfrak m$ of $S$.
\end{proof}

\noindent
Recall that we defined the
dimension $\dim_x(X)$ of a topological space $X$ at a point $x$
in Topology, Definition \ref{topology-definition-Krull}.

\begin{lemma}
\label{lemma-dimension-at-a-point-finite-type-over-field}
Let $k$ be a field.
Let $S$ be a finite type $k$-algebra.
Let $X = \Spec(S)$.
Let $\mathfrak p \subset S$ be a prime ideal and let
$x \in X$ be the corresponding point.
The following numbers are equal
\begin{enumerate}
\item $\dim_x(X)$,
\item $\max \dim(Z)$ where the maximum is over those
irreducible components $Z$ of $X$ passing through $x$, and
\item $\min \dim(S_{\mathfrak m})$ where the minimum
is over maximal ideals $\mathfrak m$ with
$\mathfrak p \subset \mathfrak m$.
\end{enumerate}
\end{lemma}

\begin{proof}
Let $X = \bigcup_{i \in I} Z_i$ be the decomposition of $X$ into
its irreducible components. There are finitely many of
them (see
Lemmas \ref{lemma-obvious-Noetherian} and \ref{lemma-Noetherian-topology}).
Let $I' = \{i \mid x \in Z_i\}$, and let
$T = \bigcup_{i \not \in I'} Z_i$. Then $U = X \setminus T$
is an open subset of $X$ containing the point $x$.
The number (2) is $\max_{i \in I'} \dim(Z_i)$.
For any open $W \subset U$ with $x \in W$
the irreducible components of $W$ are the irreducible sets
$W_i = Z_i \cap W$ for $i \in I'$ and $x$ is contained
in each of these.
Note that each $W_i$, $i \in I'$ contains a closed point because
$X$ is Jacobson, see Section \ref{section-ring-jacobson}.
Since $W_i \subset Z_i$ we have $\dim(W_i) \leq \dim(Z_i)$.
The existence of a closed point implies, via Lemma
\ref{lemma-dimension-spell-it-out}, that there is a chain of
irreducible closed subsets of length equal to $\dim(Z_i)$ in the open $W_i$.
Thus $\dim(W_i) = \dim(Z_i)$ for any $i \in I'$. Hence $\dim(W)$
is equal to the number (2). This proves that (1) $ = $ (2).

\medskip\noindent
Let $\mathfrak m \supset \mathfrak p$ be any maximal ideal
containing $\mathfrak p$. Let $x_0 \in X$ be the corresponding
point. First of all, $x_0$ is contained in all the
irreducible components $Z_i$, $i \in I'$. Let $\mathfrak q_i$
denote the minimal primes of $S$ corresponding to the
irreducible components $Z_i$. For each $i$ such that
$x_0 \in Z_i$ (which is equivalent to $\mathfrak m \supset \mathfrak q_i$)
we have a surjection
$$
S_{\mathfrak m} \longrightarrow
S_\mathfrak m/\mathfrak q_i S_\mathfrak m =(S/\mathfrak q_i)_{\mathfrak m}
$$
Moreover, the primes $\mathfrak q_i S_\mathfrak m$ so obtained
exhaust the minimal
primes of the Noetherian local ring $S_{\mathfrak m}$, see
Lemma \ref{lemma-irreducible-components-containing-x}.
We conclude, using Lemma \ref{lemma-dimension-spell-it-out},
that the dimension of $S_{\mathfrak m}$ is the
maximum of the dimensions of the $Z_i$ passing through $x_0$.
To finish the proof of the lemma it suffices to show that
we can choose $x_0$ such that $x_0 \in Z_i \Rightarrow i \in I'$.
Because $S$ is Jacobson (as we saw above)
it is enough to show that $V(\mathfrak p) \setminus T$
(with $T$ as above) is nonempty. And this is clear since it
contains the point $x$ (i.e. $\mathfrak p$).
\end{proof}

\begin{lemma}
\label{lemma-dimension-closed-point-finite-type-field}
Let $k$ be a field.
Let $S$ be a finite type $k$-algebra.
Let $X = \Spec(S)$.
Let $\mathfrak m \subset S$ be a maximal ideal and let
$x \in X$ be the associated closed point.
Then $\dim_x(X) = \dim(S_{\mathfrak m})$.
\end{lemma}

\begin{proof}
This is a special case of
Lemma \ref{lemma-dimension-at-a-point-finite-type-over-field}.
\end{proof}

\begin{lemma}
\label{lemma-disjoint-decomposition-CM-algebra}
Let $k$ be a field.
Let $S$ be a finite type $k$ algebra.
Assume that $S$ is Cohen-Macaulay.
Then $\Spec(S) = \coprod T_d$ is a finite disjoint union of
open and closed subsets $T_d$ with $T_d$ equidimensional
(see Topology, Definition \ref{topology-definition-equidimensional})
of dimension $d$. Equivalently, $S$ is a product of rings
$S_d$, $d = 0, \ldots, \dim(S)$ such that every maximal ideal
$\mathfrak m$ of $S_d$ has height $d$.
\end{lemma}

\begin{proof}
The equivalence of the two statements follows from
Lemma \ref{lemma-disjoint-implies-product}.
Let $\mathfrak m \subset S$ be a maximal ideal.
Every maximal chain of primes in $S_{\mathfrak m}$ has
the same length equal to $\dim(S_{\mathfrak m})$, see
Lemma \ref{lemma-maximal-chain-CM}. Hence, the dimension of the irreducible
components passing through the point corresponding to $\mathfrak m$
all have dimension equal to $\dim(S_{\mathfrak m})$, see
Lemma \ref{lemma-dimension-spell-it-out}.
Since $\Spec(S)$ is a Jacobson topological space
the intersection
of any two irreducible components of it contains a closed point if nonempty,
see
Lemmas \ref{lemma-finite-type-field-Jacobson} and
\ref{lemma-jacobson}.
Thus we have shown that any two irreducible components
that meet have the same dimension. The lemma follows
easily from this, and the fact that $\Spec(S)$
has a finite number of irreducible components (see
Lemmas \ref{lemma-obvious-Noetherian} and \ref{lemma-Noetherian-topology}).
\end{proof}
















